\documentclass[11pt,a4paper]{article}

\usepackage{../shared/cathedral-preamble}
\usepackage{setspace}
\onehalfspacing

\title{%
  The Prime Number Grid:\\
  How Number-Theoretic Structures from the Cathedral\\
  Could Impact Energy Systems%
}
\author{
  Jason Robert Gochanour
}
\date{April 20, 2026}

\begin{document}

\maketitle

\begin{abstract}
We examine the connections---established and speculative---between
the mathematical structures in the Cathedral project and five areas
of energy and electricity generation: (1)~power grid harmonic
analysis via M\"obius-weighted filters, (2)~turbine vibration
monitoring via arithmetic compressed sensing, (3)~formally verified
grid stability certificates, (4)~nuclear spectroscopy and the
GUE connection to zeta zeros, and (5)~fusion plasma confinement
and MHD stability via variational spectral methods. The most
immediately actionable application is harmonic decomposition in
electrical power systems, where the M\"obius function provides a
natural mathematical framework for separating fundamental-frequency
content from integer-harmonic distortion---a problem that costs the
global power industry an estimated \$30 billion annually. We
classify each proposal by feasibility, timeline, and distance from
the Cathedral's core mathematics.
\end{abstract}

\tableofcontents

% ================================================================
\section{Introduction: What Primes Have to Do with Power}
% ================================================================

At first glance, the Riemann Hypothesis has nothing to do with
electricity generation. And in the direct sense, this is correct:
proving that $d_N^2 \to 0$ will not make a turbine spin faster or a
solar panel more efficient.

But the \emph{mathematics} of the Cathedral---harmonic analysis,
spectral theory, structured matrices, variational methods, and
number-theoretic sieve constructions---is deeply connected to the
mathematics that governs energy systems. This is not a coincidence.
Energy systems are governed by differential equations whose spectral
properties (eigenvalues, resonances, stability margins) determine
performance. The Cathedral studies the spectral properties of
arithmetically structured operators. The tools overlap.

This paper maps the overlap, identifies actionable applications,
and is honest about the distance between mathematical structure
and engineering reality.

% ================================================================
\section{Power Grid Harmonics: The \$30 Billion Problem}
\label{sec:harmonics}
% ================================================================

\subsection{The Problem}

Electrical power grids operate at 50~Hz (Europe, Asia) or 60~Hz
(Americas). Nonlinear loads---variable frequency drives, LED
lighting, switching power supplies, EV chargers, solar inverters---
inject harmonic currents at integer multiples of the fundamental:
$f_k = k \cdot f_0$ for $k = 2, 3, 5, 7, 11, \ldots$

Harmonic distortion causes:
\begin{itemize}
  \item Transformer overheating (increased core and copper losses).
  \item Neutral conductor overloading (triplen harmonics: $k = 3, 9,
    15, \ldots$).
  \item Capacitor bank failures (resonance at harmonic frequencies).
  \item Protective relay misoperation.
  \item Revenue meter errors.
\end{itemize}

The global cost of power quality problems (of which harmonics are
the dominant component) is estimated at \$30 billion per year.

\subsection{Current Solutions}

\begin{itemize}
  \item \textbf{Passive filters}: LC circuits tuned to specific
    harmonics (typically 5th, 7th, 11th). Fixed, narrow-band,
    expensive.
  \item \textbf{Active power filters (APFs)}: Power electronic
    devices that inject compensating currents. Require real-time
    harmonic identification.
  \item \textbf{FFT-based identification}: Standard DFT on a
    window of voltage/current samples. Works well for stationary
    signals but struggles with time-varying harmonic content and
    inter-harmonic frequencies.
\end{itemize}

\subsection{What the Cathedral Suggests}

The harmonics in a power grid are at \emph{integer multiples} of the
fundamental frequency. This is exactly the multiplicative structure
that the M\"obius function was designed to decompose.

The M\"obius inversion formula states:
\[
  g(n) = \sum_{d | n} f(d) \quad\iff\quad
  f(n) = \sum_{d | n} \mu(n/d)\,g(d).
\]
In the power context: if $g(n)$ is the measured harmonic content at
frequency $n \cdot f_0$ (which includes contributions from all
divisors), then $f(n) = \sum_{d|n} \mu(n/d)\,g(d)$ isolates the
\emph{independent} harmonic source at frequency $n \cdot f_0$.

\textbf{Key insight}: The M\"obius function naturally separates
``fundamental'' harmonic sources from ``inherited'' harmonics. A
6th harmonic current might be a genuine 6th harmonic injection, or
it might be a product of the 2nd and 3rd harmonics interacting
through nonlinearity. M\"obius inversion distinguishes these cases.

\subsection{The Engineering Design}

\textbf{Proposal}: A M\"obius Harmonic Analyzer (MHA):

\begin{enumerate}
  \item Sample voltage and current at high resolution
    ($f_s \geq 50 \cdot f_0$).
  \item Compute the harmonic spectrum $H(k)$ for $k = 1, \ldots, K$
    via DFT.
  \item Apply M\"obius inversion: $S(k) = \sum_{d|k} \mu(k/d) H(d)$.
  \item The M\"obius-inverted spectrum $S(k)$ identifies
    \emph{independent} harmonic sources, separating directly injected
    harmonics from multiplicative products.
  \item Use $S(k)$ to drive an active power filter that targets
    source harmonics, not symptoms.
\end{enumerate}

\textbf{Advantages over FFT-based methods}:
\begin{itemize}
  \item Identifies harmonic \emph{sources} rather than symptoms.
  \item Naturally handles the multiplicative structure of harmonic
    generation.
  \item The Selberg window $(1 - \ln k / \ln K)^+$ provides
    optimal tapering for finite harmonic order.
\end{itemize}

\subsection{Feasibility}

\textbf{Very high.} This is a software algorithm that runs on
existing power quality analyzers. The M\"obius function is trivially
precomputable for $k \leq 50$ (the practical harmonic range). The
algorithm adds $O(K \log K)$ operations to the existing DFT pipeline.

\textbf{Required validation}: Compare MHA source identification with
IEEE~519 standard measurements on a real distribution feeder with
known harmonic sources (VFDs, LED banks, EV chargers).

\textbf{Potential impact}: If MHA enables more targeted active
filtering---compensating sources instead of symptoms---it could
reduce active filter sizing by 30--50\%, saving capital and energy.

% ================================================================
\section{Turbine Vibration Monitoring}
\label{sec:turbines}
% ================================================================

\subsection{The Problem}

Wind turbines, gas turbines, and steam turbines are the backbone of
electricity generation. Their mechanical health is monitored by
vibration sensors (accelerometers) that measure shaft vibration in
real time.

Vibration signals are sums of harmonics at multiples of the rotation
frequency plus defect-related sidebands. Detecting incipient faults
(bearing defects, blade cracks, misalignment) requires separating
the defect signatures from the dominant harmonic content.

Current approaches: FFT spectrum analysis, envelope analysis,
cepstral analysis. These work well for steady-state operation but
struggle during speed transients (startup, shutdown) when the
harmonic frequencies are time-varying.

\subsection{What the Cathedral Suggests}

The arithmetic compressed sensing approach from the companion paper
applies directly:

\begin{itemize}
  \item The vibration signal is naturally sparse in the
    \emph{arithmetic dictionary}: most energy is at integer multiples
    of the rotation frequency.
  \item The fractional-part basis $\{1/(kx)\}$ encodes exactly this
    multiplicative harmonic structure.
  \item The Gram matrix eigenvalue bound $\lambda_{\min}(G_N) > 0$
    guarantees that the harmonic components are always separable,
    even at low signal-to-noise ratios.
\end{itemize}

\textbf{Proposal}: An Arithmetic Vibration Monitor (AVM) that:
\begin{enumerate}
  \item Tracks the instantaneous rotation frequency $f_0(t)$ from a
    tachometer or phase-locked loop.
  \item Projects the vibration signal onto the time-varying arithmetic
    dictionary $\{1/(k \cdot f_0(t) \cdot x)\}$.
  \item Identifies anomalous components (not at integer multiples)
    as potential defect signatures.
  \item Uses M\"obius inversion to separate independent sources
    from harmonic products.
\end{enumerate}

\subsection{Feasibility}

\textbf{High.} Vibration monitoring hardware is already deployed on
most utility-scale turbines. The AVM is a software upgrade to
existing monitoring systems. The key advantage: tracking
instantaneous frequency allows analysis during transients, where
conventional FFT fails.

\textbf{Potential impact}: Earlier detection of bearing defects could
extend turbine lifetime by 10--20\% and reduce unplanned outages.
For a 5~GW wind fleet, this represents \$50--100M per year in
avoided downtime and replacement costs.

% ================================================================
\section{Formally Verified Grid Stability}
\label{sec:grid}
% ================================================================

\subsection{The Problem}

Modern power grids are increasingly complex: intermittent renewables
(wind, solar), distributed storage (batteries, EVs), demand
response, and HVDC interconnections create a system with millions of
interacting components. Grid operators must guarantee stability
under all credible operating conditions.

Current stability analysis:
\begin{itemize}
  \item \textbf{Simulation}: Run time-domain simulations for
    thousands of contingency scenarios. Computationally expensive
    and never exhaustive.
  \item \textbf{Small-signal analysis}: Linearize around the
    operating point and compute eigenvalues. Valid only near the
    linearization point.
  \item \textbf{Transient stability studies}: Offline analysis for
    planning; cannot guarantee real-time stability.
\end{itemize}

No existing approach provides a \emph{certificate} that the grid
will remain stable under a specified set of conditions.

\subsection{What the Cathedral Suggests}

The Cathedral proves positive definiteness of a family of matrices
$\{G_N\}_{N \geq 1}$---a parameterized stability certificate
verified by machine. The same framework can prove grid stability:

\begin{enumerate}
  \item \textbf{Model}: Formalize the swing equation
    $M_i \ddot{\delta}_i + D_i \dot{\delta}_i = P_i - \sum_j
    |V_i||V_j||Y_{ij}|\sin(\delta_i - \delta_j - \theta_{ij})$
    in Lean~4.
  \item \textbf{Lyapunov function}: Define the energy-like function
    $V = \frac{1}{2}\sum_i M_i \dot{\delta}_i^2 - \sum_i P_i
    (\delta_i - \delta_i^*) - \sum_{(i,j)} |V_i||V_j||Y_{ij}|
    (\cos(\delta_i - \delta_j) - \cos(\delta_i^* - \delta_j^*))$.
  \item \textbf{Prove}: $V > 0$ and $\dot{V} \leq 0$ within the
    region of attraction.
  \item \textbf{Axiomatize}: Numerical grid parameters (line
    impedances, generator inertias, load models) as named axioms
    with auditable sources.
  \item \textbf{Certify}: The output is a machine-checked statement:
    ``This grid configuration is stable under conditions A, B, C,
    conditional on axioms X, Y, Z.''
\end{enumerate}

\subsection{Feasibility}

\textbf{Medium.} The mathematical framework (Lyapunov stability for
power systems) is textbook material. The challenge is formalizing it
in Lean~4 and connecting to numerical grid data. A proof-of-concept
for a 3-bus system would be achievable in 1--2 years.

\textbf{Potential impact}: Transformative for grid operators managing
high-renewable penetration. A stability certificate could replace
thousands of simulation runs and provide provable guarantees that
no credible contingency causes instability. This is especially
valuable for islanded microgrids and ship/aircraft power systems.

% ================================================================
\section{Nuclear Spectroscopy and the GUE Connection}
\label{sec:nuclear}
% ================================================================

\subsection{The Established Science}

The Montgomery--Odlyzko law (1973/1987) established that the
pair correlation of nontrivial zeta zeros matches the GUE
(Gaussian Unitary Ensemble) statistics of random matrix theory:
\[
  R_2(\alpha) = 1 - \left(\frac{\sin(\pi\alpha)}{\pi\alpha}\right)^2
  + \delta(\alpha).
\]
The same GUE statistics describe:
\begin{itemize}
  \item Energy level spacings of heavy nuclei (confirmed
    experimentally for $^{166}$Er, $^{238}$U, and others).
  \item Eigenvalue spacings of quantum chaotic systems.
  \item Bus arrival times in Cuernavaca (a celebrated empirical
    surprise).
\end{itemize}

\subsection{What the Cathedral Adds}

The Cathedral does not work with random matrix theory directly,
but its Gram matrix $G_N$ provides a \emph{deterministic} number-
theoretic matrix whose eigenvalue distribution can be compared with
GUE predictions. The numerical experiments show:

\begin{itemize}
  \item $\lambda_{\min}(G_N) \sim C / \ln N$ (soft-edge behavior
    reminiscent of Tracy--Widom).
  \item $\lambda_{\max}(G_N) \sim N$ (linear bulk growth).
  \item The eigenvalue ratio $Q_N / \ln N \to C \approx 21.65$
    converges to a number-theoretic constant.
\end{itemize}

\subsection{Application to Nuclear Energy}

\textbf{Reactor physics}: Neutron transport in a reactor core is
governed by the Boltzmann transport equation, whose eigenvalues
determine criticality. The dominant eigenvalue $k_{\text{eff}} = 1$
at criticality; the spectral gap to the second eigenvalue determines
how quickly transients decay.

\textbf{Proposal}: Use the Cathedral's spectral analysis techniques
(Rayleigh quotient bounds, positive definiteness proofs, condition
number tracking) to provide:

\begin{enumerate}
  \item Verified bounds on the spectral gap of the neutron transport
    operator for a given core configuration.
  \item Machine-checked certificates that $k_{\text{eff}} < 1$
    (subcritical) under specified conditions.
  \item Formal proofs that shutdown margin recommendations are
    mathematically sound.
\end{enumerate}

\subsection{Feasibility}

\textbf{Low-medium for full implementation}, but the individual
mathematical components are mature. The Nuclear Regulatory
Commission (NRC) already requires validated criticality calculations;
machine-verified bounds would be a step beyond validation.

% ================================================================
\section{Fusion Plasma Confinement}
\label{sec:fusion}
% ================================================================

\subsection{The Problem}

Magnetic confinement fusion (tokamaks, stellarators) requires
maintaining a hot plasma ($10^8$ K) in a stable magnetic
configuration. The plasma is described by magnetohydrodynamics (MHD),
and stability is determined by the eigenvalues of the MHD operator.

The key challenge: proving that a given magnetic configuration is
stable against all possible plasma perturbations. This is a
variational problem---the energy principle states that the plasma
is stable if and only if:
\[
  \delta W[\xi] = \frac{1}{2} \int \left[
  |\nabla \times (\xi \times B)|^2 / \mu_0 -
  (\xi \cdot \nabla p)(\nabla \cdot \xi) + \cdots
  \right] dV > 0
\]
for all displacements $\xi$.

\subsection{What the Cathedral Suggests}

This is structurally identical to the Cathedral's core problem:
proving that a quadratic form (the MHD energy) is positive definite
for all test vectors (perturbations).

The Cathedral's techniques directly apply:
\begin{itemize}
  \item \textbf{Rayleigh--Ritz discretization}: Approximate $\xi$
    by a finite set of basis functions, reducing the continuous
    problem to a matrix eigenvalue problem $Av = \lambda Bv$.
  \item \textbf{Positive definiteness proof}: Show that
    $\lambda_{\min}(B^{-1}A) > 0$ for all relevant configurations.
  \item \textbf{Axiom tiering}: The plasma equilibrium is an axiom
    (computed numerically); the stability proof is machine-verified
    from that axiom.
\end{itemize}

\textbf{Key parallel}: The Cathedral proves $\lambda_{\min}(G_N) > 0$
for \emph{all} $N \geq 1$. An MHD stability proof would show
$\delta W > 0$ for \emph{all} perturbations within a specified class.
The mathematical structure is the same; the physics is different.

\subsection{Feasibility}

\textbf{Low for full implementation} (the MHD operator is far more
complex than the Gram matrix), but the \emph{methodology} transfers:

\begin{enumerate}
  \item Discretize the MHD stability functional on a finite element
    basis.
  \item Compute the resulting matrix and bound its minimum eigenvalue.
  \item Formalize the bound in Lean~4, with the equilibrium as an
    axiom.
  \item Produce a certificate: ``Configuration X is MHD-stable
    conditional on equilibrium axiom Y.''
\end{enumerate}

\textbf{Impact}: Machine-verified stability certificates could
accelerate the design-evaluate-redesign cycle in fusion reactor
engineering, particularly for stellarator optimization where
thousands of configurations must be evaluated.

% ================================================================
\section{Energy Harvesting with Arithmetic Metamaterials}
\label{sec:harvesting}
% ================================================================

\subsection{RF Energy Harvesting}

Radio frequency energy harvesting (rectenna arrays) captures ambient
electromagnetic energy from WiFi, cellular, broadcast, and other RF
sources. The challenge: ambient RF energy is broadband and weak,
requiring an antenna that:
\begin{itemize}
  \item Covers a wide bandwidth (100~MHz to 6~GHz).
  \item Has high efficiency at each frequency.
  \item Is physically compact.
\end{itemize}

Periodic antenna arrays have narrow bandwidth (resonant design).
Random arrays have low efficiency (poor impedance matching).

\subsection{What the Cathedral Suggests}

The prime-gap metamaterial concept from the companion paper
applies directly:

\textbf{Proposal}: A broadband rectenna with resonant elements at
prime-indexed sizes:
\begin{itemize}
  \item Element $k$ resonates at $f_k = c / (2 L_k)$ where
    $L_k \propto p_k$ (the $k$-th prime).
  \item The collective array has a dense, quasi-random frequency
    coverage: primes are distributed with density $\sim 1/\ln f$,
    providing denser coverage at low frequencies (where ambient RF
    power is higher).
  \item The natural log-density of primes \emph{matches} the
    typical ambient RF power spectrum (which decreases with
    frequency), providing more harvesting elements where there is
    more energy.
\end{itemize}

\subsection{Feasibility}

\textbf{Medium.} Wideband rectennas are an active research area.
The prime-indexed design rule is novel but physically simple to
implement. The key validation: simulate the broadband capture
efficiency versus a standard log-periodic antenna.

\textbf{Potential impact}: Prime-indexed rectennas could harvest
10--30\% more ambient RF energy than periodic designs, enabling
battery-free IoT sensors in urban environments.

% ================================================================
\section{Summary: Distance from Cathedral to Kilowatt}
\label{sec:distance}
% ================================================================

\begin{center}
\small
\begin{tabular}{@{}llccl@{}}
\toprule
\textbf{Application} & \textbf{Cathedral Math} &
\textbf{Feasibility} & \textbf{Impact} & \textbf{Timeline} \\
\midrule
Grid harmonics (MHA) & M\"obius inversion & Very High & High & $<$1 yr \\
Turbine monitoring (AVM) & Arithmetic CS & High & High & 1--2 yr \\
Grid stability certs & Verified PD proofs & Medium & Very High & 2--5 yr \\
Nuclear spectroscopy & Rayleigh quotient & Low-Med & Medium & 5+ yr \\
Fusion MHD stability & Variational PD & Low & Very High & 10+ yr \\
RF energy harvesting & Prime-gap arrays & Medium & Medium & 2--4 yr \\
\bottomrule
\end{tabular}
\end{center}

The gradient is clear: the closer the application is to harmonic
analysis (grid, turbines), the more immediately actionable. The
further into nuclear/fusion physics, the longer the timeline but
the higher the potential impact.

% ================================================================
\section{Honesty Section}
% ================================================================

\begin{enumerate}
  \item The M\"obius Harmonic Analyzer could be developed tomorrow
    without any knowledge of the Cathedral. The M\"obius function
    has been known since 1832. The Cathedral's contribution is
    organizing these ideas in a verified framework that demonstrates
    their mathematical coherence.
  \item The fusion application is a 10+ year research program that
    may ultimately prove impractical. We include it because the
    mathematical parallel (proving PD of a quadratic form = proving
    stability) is exact and illuminating.
  \item None of these proposals require \RH{} to be true. The
    mathematical structures exist independently.
  \item The economic impact estimates are order-of-magnitude only.
\end{enumerate}

% ================================================================
\section{Conclusion}
% ================================================================

The primes do not generate electricity. But the mathematics that
describes the primes---harmonic analysis, spectral theory,
variational stability, structured matrices---is the same
mathematics that governs power grids, turbines, nuclear reactors,
and fusion plasmas. The Cathedral assembles these tools in a new
configuration and verifies them by machine. This paper traces the
path from verified mathematics to possible engineering.

The most surprising connection is the simplest: power grid harmonics
occur at \emph{integer multiples} of 50/60~Hz, and the M\"obius
function is the mathematical tool built specifically for decomposing
integer-multiplicative structure. This connection has been hiding in
plain sight since August Ferdinand M\"obius defined his function in
1832 and Thomas Edison built the first power grid in 1882. It took
a formalization of the Riemann Hypothesis to make the connection
explicit.

\begin{quote}
\emph{The grid hums at 60 hertz.\\
Its harmonics are at 120, 180, 300, 420.\\
Those numbers are $2, 3, 5, 7$ times sixty.\\
The primes were in the wire all along.}
\end{quote}

\vspace{0.5cm}

\begin{thebibliography}{99}

\bibitem{ieee519}
IEEE, \emph{IEEE Std 519-2022: Recommended Practice for Harmonic
Control in Electric Power Systems}, 2022.

\bibitem{montgomery1973}
H.L.~Montgomery, \emph{The pair correlation of zeros of the zeta
function}, Proc. Symp. Pure Math. \textbf{24} (1973), 181--193.

\bibitem{mehta2004}
M.L.~Mehta, \emph{Random Matrices}, 3rd ed., Academic Press, 2004.

\bibitem{freidberg2014}
J.P.~Freidberg, \emph{Ideal MHD}, Cambridge, 2014.

\bibitem{kundur1994}
P.~Kundur, \emph{Power System Stability and Control}, McGraw-Hill,
1994.

\bibitem{cathedral-engineering}
J.R.~Gochanour, \emph{From Zeta Zeros to Engineering}, 2026.

\bibitem{cathedral-futures}
J.R.~Gochanour, \emph{What Could Be Built: Seven Engineering
Frontiers}, 2026.

\end{thebibliography}

\end{document}
